\documentclass[a4paper,10pt]{article}
\usepackage[utf8]{inputenc}
\usepackage[T1]{fontenc}
\usepackage[french]{babel}
\usepackage[margin=1.5cm]{geometry}
\usepackage{xcolor}
\usepackage{tcolorbox}
\usepackage{enumitem}
\usepackage{multicol}
\usepackage{lmodern}
\usepackage{tabularx}
\usepackage{booktabs}

\renewcommand{\familydefault}{\sfdefault}
\pagestyle{empty}
\setlist{nosep, leftmargin=1.5em}

\definecolor{navy}{RGB}{0,51,102}
\definecolor{teal}{RGB}{0,120,120}
\definecolor{crimson}{RGB}{180,0,0}
\definecolor{green}{RGB}{0,120,50}
\definecolor{lightbg}{RGB}{245,248,252}
\definecolor{orange}{RGB}{200,100,0}
\definecolor{purple}{RGB}{100,0,160}
\definecolor{darkblue}{RGB}{0,60,130}

\tcbset{basebox/.style={
  colback=lightbg, colframe=darkblue, coltitle=white,
  fonttitle=\bfseries, arc=2mm, boxrule=0.5mm,
  left=3mm, right=3mm, top=2mm, bottom=2mm, after skip=6pt
}}

\newcommand{\key}[1]{\textbf{\textcolor{darkblue}{#1}}}
\newcommand{\warn}[1]{\textbf{\textcolor{crimson}{#1}}}
\newcommand{\ok}[1]{\textbf{\textcolor{green}{#1}}}

\begin{document}

\begin{center}
  \colorbox{darkblue}{\parbox{\dimexpr\textwidth-2\fboxsep}{
    \centering\vspace{3pt}
    {\LARGE\bfseries\textcolor{white}{CEJM — Contexte Économique, Juridique et Managérial}}\\
    {\normalsize\textcolor{white}{BTS SIO — Maël Mignard}}
    \vspace{3pt}
  }}
\end{center}
\vspace{6pt}

\begin{multicols}{2}

% ============================================================
\begin{tcolorbox}[basebox, colframe=darkblue, title=1. ÉCONOMIE — Les agents économiques]
\begin{tabular}{lll}
\toprule
\textbf{Agent} & \textbf{Fonction} & \textbf{Ressource} \\
\midrule
Ménages & Consommation & Salaires \\
Entreprises & Production marchande & Profits \\
État & Production non marchande & Impôts \\
Banques & Financement & Intérêts \\
Reste du monde & Échanges ext. & Export/Import \\
\bottomrule
\end{tabular}
\vspace{4pt}\\
\key{Flux réels} : biens et services \\
\key{Flux monétaires} : paiements correspondants \\
\key{Circuit économique} : interactions entre agents
\end{tcolorbox}

% ============================================================
\begin{tcolorbox}[basebox, colframe=darkblue, title=2. ÉCONOMIE — Marchés et prix]
\key{Marché} : lieu (physique ou virtuel) de rencontre entre l'offre et la demande.\\[4pt]
\textbf{Loi de l'offre et la demande :}
\begin{itemize}
  \item Demande $>$ Offre $\to$ Prix $\nearrow$
  \item Offre $>$ Demande $\to$ Prix $\searrow$
  \item Équilibre : Offre = Demande $\to$ Prix d'équilibre
\end{itemize}
\vspace{4pt}
\textbf{Structures de marché :}
\begin{itemize}
  \item \key{Concurrence pure et parfaite} : nombreux acteurs, produits identiques
  \item \key{Monopole} : 1 seul vendeur (ex: SNCF sur certaines lignes)
  \item \key{Oligopole} : quelques vendeurs (ex: opérateurs télécom)
  \item \key{Duopole} : 2 vendeurs (ex: Airbus/Boeing)
\end{itemize}
\vspace{4pt}
\key{Inflation} : hausse générale et durable des prix \\
\key{PIB} : Produit Intérieur Brut, richesse créée par un pays
\end{tcolorbox}

% ============================================================
\begin{tcolorbox}[basebox, colframe=teal, title=3. ÉCONOMIE — Mondialisation et numérique]
\key{Mondialisation} : interdépendance économique mondiale (libre-échange, division internationale du travail) \\[4pt]
\key{Commerce international} : exportations et importations \\
\key{Balance commerciale} : exports - imports \\
\quad $> 0$ : excédent ; $< 0$ : déficit \\[4pt]
\key{Économie numérique} :
\begin{itemize}
  \item \key{Plateformes} : Uber, Amazon, Airbnb (effet de réseau)
  \item \key{Big Data} : exploitation des données massives
  \item \key{IA} : automatisation, nouveaux métiers
  \item \key{Uberisation} : mise en relation directe via appli
\end{itemize}
\vspace{4pt}
\key{GAFAM} : Google, Apple, Facebook, Amazon, Microsoft
\end{tcolorbox}

% ============================================================
\begin{tcolorbox}[basebox, colframe=crimson, title=4. DROIT — Sources du droit]
\textbf{Hiérarchie des normes (pyramide de Kelsen) :}
\begin{enumerate}
  \item Constitution (loi fondamentale)
  \item Traités internationaux
  \item Lois (Parlement) et ordonnances
  \item Décrets (gouvernement)
  \item Arrêtés (ministres/préfets)
\end{enumerate}
\vspace{4pt}
\key{Droit public} : relations État/citoyens \\
\key{Droit privé} : relations entre particuliers \\
\quad - Droit civil, droit commercial, droit du travail \\[4pt]
\key{Droit pénal} : infractions et sanctions (contravention, délit, crime)
\end{tcolorbox}

% ============================================================
\begin{tcolorbox}[basebox, colframe=crimson, title=5. DROIT — Formation du contrat]
\textbf{Conditions de validité du contrat (art. 1128 Code civil) :}
\begin{itemize}
  \item \key{Consentement} : libre, éclairé, non vicié
  \item \key{Capacité} : être majeur ou émancipé
  \item \key{Objet} : licite et déterminé
  \item \key{Cause} : licite et réelle
\end{itemize}
\vspace{4pt}
\textbf{Vices du consentement (rendent le contrat annulable) :}
\begin{itemize}
  \item \warn{Erreur} : mauvaise compréhension de l'objet
  \item \warn{Dol} : tromperie, manœuvres frauduleuses
  \item \warn{Violence} : contrainte physique ou morale
\end{itemize}
\vspace{4pt}
\key{Nullité relative} : peut être demandée par la victime \\
\key{Nullité absolue} : si ordre public violé (annulation par tous)
\end{tcolorbox}

% ============================================================
\begin{tcolorbox}[basebox, colframe=crimson, title=6. DROIT — Droit du travail]
\key{CDI} : Contrat à Durée Indéterminée (droit commun) \\
\key{CDD} : Contrat à Durée Déterminée (cas limitatifs) \\
\key{Stage} : non salarié, en formation \\
\key{Intérim} : via agence, entreprise utilisatrice \\[4pt]
\textbf{Rupture du contrat de travail :}
\begin{itemize}
  \item \key{Démission} : à l'initiative du salarié
  \item \key{Licenciement} : à l'initiative de l'employeur
    \begin{itemize}
      \item Personnel (motif réel et sérieux)
      \item Économique (raisons économiques)
    \end{itemize}
  \item \key{Rupture conventionnelle} : accord commun
\end{itemize}
\vspace{4pt}
\key{Préavis} : délai de prévenance avant départ \\
\key{Indemnités} : compensations financières \\
\key{Conventions collectives} : accords de branche (supérieurs au Code du travail)
\end{tcolorbox}

% ============================================================
\begin{tcolorbox}[basebox, colframe=green, title=7. MANAGEMENT — Entreprise et formes juridiques]
\begin{tabular}{lll}
\toprule
\textbf{Forme} & \textbf{Responsabilité} & \textbf{Capital min.} \\
\midrule
EI & Illimitée & Aucun \\
EURL & Limitée & 1€ \\
SARL & Limitée & 1€ \\
SA & Limitée & 37 000€ \\
SAS/SASU & Limitée & 1€ \\
\bottomrule
\end{tabular}
\vspace{4pt}\\
\key{Entrepreneur individuel} : pas de séparation patrimoine \\
\key{SARL} : 2 à 100 associés, gérant, parts sociales \\
\key{SA} : actionnaires, PDG, Conseil d'administration \\
\key{Micro-entreprise} : régime simplifié (ex auto-entrepreneur)
\end{tcolorbox}

% ============================================================
\begin{tcolorbox}[basebox, colframe=green, title=8. MANAGEMENT — Finalité et performance]
\key{Finalité lucrative} : maximiser le profit \\
\key{Finalité non lucrative} : association, syndicat, service public \\[4pt]
\textbf{Types de performance :}
\begin{itemize}
  \item \key{Efficacité} : atteindre les objectifs fixés
  \item \key{Efficience} : atteindre les objectifs avec le minimum de ressources
  \item \key{RSE} : performance sociale et environnementale
\end{itemize}
\vspace{4pt}
\key{Indicateurs} : chiffre d'affaires, part de marché, taux de satisfaction, ROI \\[4pt]
\key{SWOT} : Forces, Faiblesses, Opportunités, Menaces
\end{tcolorbox}

% ============================================================
\begin{tcolorbox}[basebox, colframe=orange, title=9. MANAGEMENT — Styles et organisation]
\textbf{Styles de management (Likert) :}
\begin{tabular}{ll}
\toprule
\textbf{Style} & \textbf{Description} \\
\midrule
Directif & Ordres, peu d'autonomie, contrôle fort \\
Persuasif & Explique, motive, guide \\
Participatif & Concertation, équipe impliquée \\
Délégatif & Forte autonomie, confiance totale \\
\bottomrule
\end{tabular}
\vspace{4pt}\\
\key{Organisation} = structure + coordination \\
\textbf{Structures :}
\begin{itemize}
  \item \key{Hiérarchique} : autorité verticale (pyramide)
  \item \key{Fonctionnelle} : spécialisation par fonctions
  \item \key{Matricielle} : projets + fonctions croisés
\end{itemize}
\end{tcolorbox}

% ============================================================
\begin{tcolorbox}[basebox, colframe=purple, title=10. MANAGEMENT — Parties prenantes et RSE]
\key{Parties prenantes (stakeholders)} : tous ceux impactés par l'entreprise
\begin{itemize}
  \item \textbf{Internes} : dirigeants, salariés, actionnaires
  \item \textbf{Externes} : clients, fournisseurs, État, société civile
\end{itemize}
\vspace{4pt}
\key{RSE} (Responsabilité Sociétale des Entreprises) :\\
Intégrer les enjeux sociaux et environnementaux dans la stratégie.
\begin{itemize}
  \item Environnement : réduire l'empreinte carbone
  \item Social : conditions de travail, diversité
  \item Gouvernance : éthique, transparence
\end{itemize}
\vspace{4pt}
\key{Développement durable} : satisfaire les besoins du présent sans compromettre ceux du futur
\end{tcolorbox}

\end{multicols}

\vspace{4pt}
\begin{tcolorbox}[basebox, colframe=darkblue, title=Points clés à retenir pour l'examen CEJM]
\begin{multicols}{2}
\begin{itemize}
  \item Agents éco : Ménages, Entreprises, État, Banques, Reste du monde
  \item Offre $>$ Demande $\to$ Prix baisse ; Demande $>$ Offre $\to$ Prix monte
  \item Monopole = 1 vendeur ; Oligopole = peu ; CPP = beaucoup
  \item Pyramide des normes : Constitution $>$ Loi $>$ Décret $>$ Arrêté
  \item Contrat valide : consentement + capacité + objet licite + cause licite
  \item Vices : erreur, dol (tromperie), violence
  \item CDI = droit commun ; CDD = cas exceptionnels limitatifs
  \item Rupture conventionnelle = commun accord salarié + employeur
  \item Efficacité $\neq$ Efficience : atteindre l'objectif $\neq$ le faire avec le minimum
  \item RSE = Triple bottom line : économique + social + environnemental
  \item SARL : 1€ de capital, responsabilité limitée, parts sociales
  \item Styles management : Directif / Persuasif / Participatif / Délégatif
\end{itemize}
\end{multicols}
\end{tcolorbox}

\end{document}
